\chapter{مدیریت مراجع و منابع تحقیق}
برای بسیاری از دانشجویان ارجاع دهی و فهرست مراجع استاندارد یكی از وقت گیرترین قسمت ها در نگارش یك پایان نامه
است و از طرف دیگر به گفته داوران، این بخش یكی از رای جترین قسمت هایی است كه دارای خطای ویرایشی است. امروزه
نرم افزارهای بسیاری برای مدیریت منابع مورد استفاده قرار می گیرد. در نرم‌افزار لاتک مراجع  را می‌توان به صورت دستی تنظیم نمود. با این حال برای متون طولانی نظیر پایان‌نامه توصیه می‌شود از قابلیت \lr{bibtex} استفاده گردد. برای مثال مراجع این متن در یک فایل با نام MyReferences.bib قرار داده شده‌اند. این فایل می‌تواند شامل تعداد زیادی مرجع باشد اما با استفاده از قابلیت \lr{bibtex} تنها مراجعی نشان داده می‌شوند که در متن به آنها ارجاع داده شده باشد. شایان ذکر است اغلب پایگاه داده‌های منابع علمی نظیر \lr{scopus} و \lr{science direct} امکان دریافت مراجع در قالب \lr{bibtex} وجود دارد.

سبک‌های مختلفی برای نحوه نمایش، شماره‌گذاری و مرتب کردن مراجع وجود دارد. این سبک‌ها به صورت فایل‌هایی با پسوند bst ارائه می‌شوند. با این حال جهت هماهنگی با متون فارسی توصیه می‌شود از سبک‌هایی که در آنها تمهیدات لازم برای مراجع فارسی پیش‌بینی شده است، استفاده گردد. در سایت پارسی لاتک چندین سبک برای ارجاع‌دهی همزمان به اسناد فارسی و انگلیسی ارائه گردیده است. در ادامه نمونه‌ای از ارجاع به اسناد مختلف آورده شده است. 

مرجع \cite{Omidali82phdThesis} یک نمونه پروژه دکترا و مرجع \cite{Vahedi87} یک نمونه مقاله مجله فارسی است.
مرجع \cite{Amintoosi87afzayesh}  یک نمونه  مقاله کنفرانس فارسی و
مرجع \cite{Pedram80osool} یک نمونه کتاب فارسی با ذکر مترجمان و ویراستاران فارسی است. مرجع \cite{Khalighi07MscThesis} یک نمونه پروژه کارشناسی ارشد انگلیسی و
\cite{Khalighi87xepersian} هم یک نمونه متفرقه  می‌باشند.

مرجع \cite{Gonzalez02book} یک نمونه کتاب لاتین است که از آنجا که دارای فیلد \lr{authorfa} است، نام نویسندگان آن در استیلهای \lr{asa-fa}، \lr{plainnat-fa} و \lr{chicago-fa} به فارسی دیده می‌شود. مرجع \cite{Baker02limits} مقاله انگلیسی است که معادل فارسی نام نویسندگان آن ذکر نشده بوده است.


{\small
% در اینجا می‌توانید سبک‌های مختلف را گذاشته و تفاوت خروجی را ببینید
\bibliographystyle{acm-fa}

\bibliography{MyReferences}
}
